\chapter{Acidi e basi}
 
La prima definizione quantitativa di acidi e basi è quella di Arrhenius\mkeyword{teoria di Arrhenius}, secondo la quale un acido è una sostanza che aumenta la concentrazione di $\mathrm{H^+}$ in soluzione e una base una sostanza che aumenta quella di $\mathrm{OH^-}$, cosicché la reazione fra i due produce un sale e acqua.
Il limite di questa impostazione è che vincola l'intera chimica acido-base alle soluzioni acquose e richiede che ogni base contenga il gruppo $\mathrm{OH}$.\\
 
La teoria di Br\o nsted-Lowry\mkeyword{teoria di Br\o nsted-Lowry} generalizza il concetto definendo una reazione acido-base come un semplice trasferimento di protoni: acido è la specie che cede $\mathrm{H^+}$, base quella che lo accetta.
Ne segue che una base non deve necessariamente contenere ossidrili, e che la stessa sostanza può comportarsi da acido o da base a seconda del partner.
L'acqua ne è l'esempio principale, poiché di fronte a $\mathrm{HCl}$ accetta un protone mentre di fronte a $\mathrm{NH_3}$ lo cede:

\begin{equation}
    \mathrm{HCl} + \mathrm{H_2O} \longrightarrow \mathrm{H_3O^+} + \mathrm{Cl^-} ,
    \qquad
    \mathrm{NH_3} + \mathrm{H_2O} \rightleftharpoons \mathrm{NH_4^+} + \mathrm{OH^-} .
    \label{eq:brønsted_esempi}
\end{equation}

Una sostanza con questa doppia attitudine si dice anfiprotica\mkeyword{sostanza anfiprotica}, e nel caso dell'acqua la conseguenza è che essa reagisce con se stessa secondo l'equilibrio di autoionizzazione che studieremo più avanti.\\
 
Si osservi che il protone non esiste mai libero in soluzione: essendo un nucleo nudo, si lega immediatamente al doppietto libero di una molecola d'acqua formando lo ione ossonio $\mathrm{H_3O^+}$.
Le notazioni $\mathrm{H^+}$ e $\mathrm{H_3O^+}$ si usano indifferentemente, con la seconda più aderente alla realtà.\\
 
Ogni reazione di trasferimento protonico mette in gioco due coppie acido-base coniugate\mkeyword{coppie coniugate}: ciò che resta dell'acido dopo la cessione del protone è la sua base coniugata, e ciò che si forma dalla base dopo l'acquisto è il suo acido coniugato.
Vale una regola di reciprocità che conviene tenere presente: quanto più forte è l'acido, tanto più debole è la sua base coniugata, e viceversa.
Ne segue che alle basi coniugate degli acidi forti non resta alcuna tendenza ad accettare protoni, mentre le coppie coniugate delle specie deboli sono a loro volta deboli.
 
\section{Acidi e basi forti}
 
Gli acidi forti sono elettroliti forti, cioè in acqua esistono soltanto come ioni, e la quantità di acido indissociato all'equilibrio è trascurabile.
Appartengono a questa categoria $\mathrm{HBr}$, $\mathrm{HI}$ e $\mathrm{HClO_3}$, e in generale quasi tutti gli idracidi.
 
Le basi forti tipiche sono gli idrossidi dei metalli alcalini e alcalino-terrosi, come $\mathrm{NaOH}$, $\mathrm{KOH}$ e $\mathrm{Ca(OH)_2}$, che sono elettroliti forti e liberano $\mathrm{OH^-}$ per semplice dissociazione, purché siano solubili.

Esistono però anche specie che sono basi forti in un senso più radicale, cioè più basiche dello stesso ione ossidrile, come lo ione ossido, lo ione idruro e lo ione ammiduro; queste reagiscono violentemente con l'acqua, in modo fortemente esotermico, secondo reazioni del tipo
\begin{equation}
    \mathrm{O^{2-}} + \mathrm{H_2O} \longrightarrow 2\,\mathrm{OH^-} ,
    \qquad
    \mathrm{H^-} + \mathrm{H_2O} \longrightarrow \mathrm{H_2} + \mathrm{OH^-} ,
    \label{eq:basi_fortissime}
\end{equation}
 
\section{Acidi e basi deboli}
 
Un acido debole è solo parzialmente dissociato, cosicché in soluzione coesistono molecole dissociate e indissociate in un equilibrio vero e proprio,

\begin{equation}
    \mathrm{HA} + \mathrm{H_2O} \rightleftharpoons \mathrm{A^-} + \mathrm{H_3O^+} ,
    \qquad
    K_a = \frac{[\mathrm{H_3O^+}][\mathrm{A^-}]}{[\mathrm{HA}]} ,
    \label{eq:ka}
\end{equation}

dove la costante $K_a$ si dice costante di dissociazione acida\mkeyword{costante di dissociazione acida} e l'acqua non compare perché solvente in largo eccesso.
Simmetricamente una base debole sottrae parzialmente protoni all'acqua,

\begin{equation}
    \mathrm{B} + \mathrm{H_2O} \rightleftharpoons \mathrm{BH^+} + \mathrm{OH^-} ,
    \qquad
    K_b = \frac{[\mathrm{BH^+}][\mathrm{OH^-}]}{[\mathrm{B}]} .
    \label{eq:kb}
\end{equation}
 
Le due costanti di una stessa coppia coniugata non sono indipendenti.
Sommando membro a membro l'equilibrio di dissociazione dell'acido e quello della sua base coniugata si ottiene infatti la sola autoionizzazione dell'acqua, e poiché alla somma di due reazioni corrisponde il prodotto delle costanti,

\begin{equation}
    K_a \cdot K_b = K_w .
    \label{eq:kakb}
\end{equation}

La \eqref{eq:kakb} è la forma quantitativa della reciprocità enunciata più sopra, e permette di ricavare la costante basica di un anione conoscendo quella acida del suo acido coniugato.
 
\section{Fattori che determinano la forza acida}
 
Perché una specie $\mathrm{HX}$ sia un acido occorrono tre condizioni: 

\begin{itemize}
    \item che il $\mathrm{H}$--$\mathrm{X}$ sia polarizzato con carica parziale positiva sull'idrogeno, il che richiede $\mathrm{X}$ più elettronegativo di $\mathrm{H}$;
    \item che il legame sia abbastanza debole da poter essere rotto;
    \item che la base coniugata $\mathrm{X^-}$ risultante sia stabile.
\end{itemize}

Le tre condizioni non vanno sempre nella stessa direzione, e la forza acida osservata è il risultato del loro compromesso.\\
 
Negli idracidi\mkeyword{idracidi} la forza acida aumenta lungo un periodo, dove cresce l'elettronegatività e con essa la polarizzazione del legame, e aumenta scendendo lungo un gruppo, dove cresce il volume atomico e il legame si indebolisce.
Il secondo andamento prevale sul primo, ed è la ragione per cui $\mathrm{HF}$, pur essendo il più polarizzato della serie, è un acido debole mentre $\mathrm{HCl}$, $\mathrm{HBr}$ e $\mathrm{HI}$ sono forti: il legame $\mathrm{H}$--$\mathrm{F}$ è troppo forte per essere rotto facilmente.
All'estremo opposto il metano è neutro, poiché il legame $\mathrm{C}$--$\mathrm{H}$ non è praticamente polarizzato.

Gli ossiacidi\mkeyword{ossiacidi} hanno tutti struttura generale $\mathrm{Y}$--$\mathrm{O}$--$\mathrm{H}$, e l'unico legame che si rompe è quello $\mathrm{O}$--$\mathrm{H}$; la forza dipende quindi da $\mathrm{Y}$ e dagli atomi a esso legati.
Se $\mathrm{Y}$ è un metallo, poco elettronegativo, a rompersi è invece il legame $\mathrm{Y}$--$\mathrm{O}$ e il composto è una base; se $\mathrm{Y}$ ha elettronegatività intermedia il composto è un acido debole; se $\mathrm{Y}$ è molto elettronegativo sottrae densità elettronica all'ossigeno, che a sua volta la sottrae all'idrogeno, e il composto è un acido forte.
Aumentando il numero di ossigeni legati a $\mathrm{Y}$ la forza acida cresce ulteriormente, di circa cinque ordini di grandezza per ogni ossigeno aggiunto, sia perché aumenta la polarità del legame $\mathrm{O}$--$\mathrm{H}$ sia, soprattutto, perché la carica negativa della base coniugata può delocalizzarsi su un numero maggiore di ossigeni equivalenti.
Ne seguono i due andamenti $\mathrm{H_3PO_4} < \mathrm{H_2SO_4} < \mathrm{HClO_4}$ lungo il periodo e $\mathrm{HClO_4} > \mathrm{HBrO_4} > \mathrm{HIO_4}$ lungo il gruppo.
 

\section{Equilibri acido-base in acqua}
 
L'acqua, essendo anfiprotica, reagisce con se stessa secondo l'equilibrio di autoionizzazione, la cui costante non contiene il solvente e prende il nome di prodotto ionico dell'acqua\mkeyword{prodotto ionico dell'acqua}:
\begin{equation}
    2\,\mathrm{H_2O} \rightleftharpoons \mathrm{H_3O^+} + \mathrm{OH^-} ,
    \qquad
    K_w = [\mathrm{H_3O^+}][\mathrm{OH^-}] = 1{,}0\cdot10^{-14} \ \text{a } 25 \ ^\circ C .
    \label{eq:kw}
\end{equation}
La \eqref{eq:kw} vale in qualunque soluzione acquosa, e implica che le concentrazioni dei due ioni siano inversamente proporzionali: in soluzione non possono mai coesistere in quantità entrambe rilevanti.
 
Sciogliendo un acido debole $\mathrm{HA}$ a concentrazione formale $C_A$ si hanno due equilibri simultanei, quello dell'acido e quello dell'acqua, e quattro specie incognite, $[\mathrm{H^+}]$, $[\mathrm{OH^-}]$, $[\mathrm{HA}]$ e $[\mathrm{A^-}]$.
Poiché due equazioni non bastano a determinare quattro incognite, si aggiungono due condizioni di conservazione: il bilancio di massa, per cui l'acido introdotto si ritrova nelle due forme
\begin{equation}
    C_A = [\mathrm{HA}] + [\mathrm{A^-}] ,
    \label{eq:bilancio_massa}
\end{equation}
e il bilancio di carica, per cui la soluzione è complessivamente neutra,
\begin{equation}
    [\mathrm{H^+}] = [\mathrm{A^-}] + [\mathrm{OH^-}] .
    \label{eq:bilancio_carica}
\end{equation}
Quest'ultimo si legge anche come bilancio dei protoni, poiché gli ioni $\mathrm{H^+}$ provengono sia dall'acido, lasciando $\mathrm{A^-}$, sia dall'acqua, lasciando $\mathrm{OH^-}$.
 
Il sistema di quattro equazioni si riduce, eliminando le altre incognite, a un'equazione di terzo grado in $[\mathrm{H^+}]$,
\begin{equation}
    [\mathrm{H^+}]^3 + K_a[\mathrm{H^+}]^2 - (K_aC_A + K_w)[\mathrm{H^+}] - K_aK_w = 0 ,
    \label{eq:cubica}
\end{equation}
che non si risolve a mano ma che ammette due approssimazioni successive.
La prima consiste nel trascurare l'autoionizzazione dell'acqua, lecita quando $K_aC_A \gg K_w$, e riduce la \eqref{eq:cubica} a un'equazione di secondo grado; la seconda consiste nel trascurare la frazione di acido dissociato rispetto a $C_A$, lecita quando l'acido è debole e non troppo diluito, e conduce all'espressione di uso corrente
\begin{equation}
    [\mathrm{H^+}] = \sqrt{K_a C_A} .
    \label{eq:acido_debole}
\end{equation}
Per un acido forte, viceversa, la dissociazione è completa e si ha semplicemente $[\mathrm{H^+}] = C_A$.
 
\section{La scala del pH}
 
Poiché nella maggior parte delle soluzioni le concentrazioni ioniche sono piccole e si estendono su molti ordini di grandezza, conviene usare una scala logaritmica:

\begin{equation}
    pH = -\log[\mathrm{H^+}] ,
    \qquad
    pOH = -\log[\mathrm{OH^-}] ,
    \qquad
    pH + pOH = 14{,}00 .
    \label{eq:ph}
\end{equation}

L'ultima relazione discende direttamente dal prodotto ionico, e mostra che le due grandezze non sono indipendenti.
A $25 \ ^\circ C$ una soluzione neutra ha $pH = pOH = 7,00$, una acida ha $pH < 7$ e una basica $pH > 7$.
 
Applicando la \eqref{eq:ph} ai due casi limite si ottengono le formule di uso quotidiano: per un acido forte $pH = -\log C_A$, e per un acido debole, dalla \eqref{eq:acido_debole},

\begin{equation}
    pH = \tfrac{1}{2}\left(pK_a - \log C_A\right) .
    \label{eq:ph_debole}
\end{equation}

Si noti la differenza strutturale fra le due: nel primo caso il pH dipende dalla sola concentrazione, nel secondo anche dalla natura dell'acido.\\

Sperimentalmente il pH si misura con un pH-metro, costituito da un elettrodo a vetro il cui potenziale varia in funzione della concentrazione di $\mathrm{H^+}$, oppure, in modo più grossolano, con indicatori\mkeyword{indicatori}, cioè sostanze che cambiano colore in un intervallo di una o due unità di pH.

\section{Le soluzioni tampone}
 
Una soluzione tampone\mkeyword{soluzione tampone} è una miscela di un acido debole e della sua base coniugata in quantità confrontabili, ottenuta di norma sciogliendo l'acido insieme a un suo sale con base forte.
La sua proprietà caratteristica è che il pH risulta poco sensibile all'aggiunta di acidi o basi, poiché la perturbazione viene neutralizzata da una delle due componenti e si traduce in una variazione del loro rapporto anziché in ioni liberi.
 
Riscrivendo la costante di dissociazione ed esplicitando la concentrazione di protoni si ottiene
\begin{equation}
    [\mathrm{H^+}] = K_a\,\frac{C_a}{C_s}
    \qquad\Longrightarrow\qquad
    pH = pK_a + \log\frac{C_s}{C_a} ,
    \label{eq:henderson}
\end{equation}
dove $C_a$ e $C_s$ sono le concentrazioni dell'acido e del suo sale.
La \eqref{eq:henderson} mette in evidenza le due proprietà fondamentali di un tampone: il pH dipende dal solo rapporto fra le due specie, e quindi non cambia per diluizione, e vale esattamente $pK_a$ quando le due specie sono equimolari.
Ne segue il criterio con cui si progetta un tampone, cioè scegliere l'acido il cui $pK_a$ sia più vicino al pH desiderato, e l'intervallo entro il quale esso è efficace, dell'ordine di $pK_a \pm 1$.
 
La capacità tamponante non è però illimitata: aggiungendo base si consuma progressivamente l'acido, e quando esso è esaurito resta soltanto il sale, la cui idrolisi rende la soluzione basica e da quel punto in poi il pH varia bruscamente.
Un ragionamento del tutto simmetrico vale per il tampone basico, costituito da una base debole e dal suo acido coniugato, per il quale si ha $[\mathrm{OH^-}] = K_b\,C_b/C_s$.